%%%%%%%%%%%%%%%%%%%%%%%
%%% DOCUMENT SETTINGS
%%%%%%%%%%%%%%%%%%%%%%%
\documentclass[11pt]{exam}
\printanswers % Uncomment to show answers
%%%%%%%%%%%%%%
%%% PACKAGES
%%%%%%%%%%%%%%
\usepackage{amsmath,amsfonts,amssymb,amsthm} % math fonts and symbols
\usepackage{bbm} % Math blackboard font
\usepackage{enumitem} % For reference enumerations lists
\usepackage{textcomp} % some symbols (like copyright)
\usepackage[top = 2cm, bottom = 2cm, right = 1.2cm, left = 1.2cm, paper = a4paper, headheight = 6cm]{geometry} % Margins setup
\usepackage{xcolor} % Colors
\usepackage{graphicx}
%%%%%%%%%%%%%%%%%%%%%%%%%%
%%% MACROS AND COMMANDS
%%%%%%%%%%%%%%%%%%%%%%%%%%
\newcommand{\N}{\mathbb{N}} % natural number set
\newcommand{\R}{\mathbb{R}} % real number set
\makeatletter
\renewcommand{\@makefnmark}{\textsuperscript{\arabic{footnote}}} % Ensure numeric superscripts
\renewcommand{\thefootnote}{\arabic{footnote}} % Ensure numeric labels
\makeatother
\setcounter{footnote}{1}
\newcommand{\BAR}{%
    \hspace{-\arraycolsep}%
    \strut\vrule % the `\vrule` is as high and deep as a strut
    \hspace{-\arraycolsep}%
}
%%%%%%%%%%%%%%%%%%%%%%%%%%%%%%%%%
%%% HEADER AND FOOT 
%%%%%%%%%%%%%%%%%%%%%%%%%%%%%%%%%
% Defining information
\newcommand{\degree}{Bachelor in Philosophy, Politics and Economics}
\newcommand{\shortdeg}{PPE}
\newcommand{\exam}{1st Midterm}
\newcommand{\subject}{Mathematics}
\newcommand{\theyear}{2024/25}
\newcommand{\ccauthor}{A. G. Azze}
\newcommand{\thedate}{Oct 30, 2024}
\newcommand{\university}{CUNEF Universidad}
% Setting header and footer
\pagestyle{headandfoot}
\header{\degree \\ \subject\ - \theyear}{}{\thedate \\ \ccauthor}
\footer{\university}{}{\thepage}
%%%%%%%%%%%%%%%%%%%%%%%%%
%%% DOCUMENT CONTENT
%%%%%%%%%%%%%%%%%%%%%%%%%
\begin{document}
%--- Instructions ---%
\begin{center}
    {\Huge \textbf{\exam{}}}
    \vspace{0.7cm}
    \begin{flushright}
        \textbf{Time Limit:} 50 minutes
    \end{flushright}
    \vspace{-0.2cm}
    
    \fbox{\parbox{\textwidth}{\centering 
        \begin{enumerate}[leftmargin = 0.5cm, label = $\bullet$, rightmargin = 0.2cm]
            \item \textbf{Every non-trivial calculation and claim in the exam must be properly justified.}

            \item The use of calculators is not necessary and therefore not allowed.

            \item Digital devices such as mobile phones, tablets, and laptops, as well as internet access, are forbidden.
        \end{enumerate}
    }}
\end{center}
%--- Instructions ---%

%--- Questions ---%

\begin{questions}

% Question 1
\question[3]
    Consider the function $f(x) = \sqrt{x - 5}\ln(x - 5)$.
    \begin{enumerate}[leftmargin = 0.5cm, label = (\alph*), rightmargin = 0.2cm]
        \item Find the domain of $f$.
        \item Find the continuity domain of $f$.
        \item Calculate the limit $\lim_{x\rightarrow 5} f(x)$. (Hint: Note that $\sqrt{x - 5}\ln(x - 5) = \frac{\ln(x - 5)}{(x - 5)^{-1/2}}$ and use the L'Hospital rule)
        \item For $h(x) = x^3 + 1$ and $g(x) = 2x$, find the limit $\lim_{x\rightarrow 5} (h \circ g \circ f)(x)$.
    \end{enumerate}

% Solution
\begin{solution}[0cm]
    \begin{enumerate}[leftmargin = 0.5cm, label = (\alph*), rightmargin = 0.2cm]
        \item The domain of $f$ is $D = (5, \infty)$.
        \item The continuity domain is also $D$.
        \item Using the L'Hospital rule, we find $\lim_{x\rightarrow 5} f(x) = 0$.
        \item $\lim_{x\rightarrow 5} (h \circ g \circ f)(x) = h(0) = 1$.
    \end{enumerate}
\end{solution}

% Question 2
\question[3]
    Let $f(x) = xe^{-x}$.
    \begin{enumerate}[leftmargin = 0.5cm, label = (\alph*), rightmargin = 0.2cm]
        \item Calculate $f'(x)$.
        \item Use $f'(0)$ to approximate $f(0.01)$.
        \item Evaluate $\lim_{x \rightarrow \infty} f(x)$.
    \end{enumerate}

% Solution
\begin{solution}[0cm]
    \begin{enumerate}[leftmargin = 0.5cm, label = (\alph*), rightmargin = 0.2cm]
        \item $f'(x) = e^{-x}(1 - x)$.
        \item $f(0.01) \approx f(0) + f'(0)(0.01 - 0) = 0.01$.
        \item $\lim_{x \to \infty} f(x) = 0$.
    \end{enumerate}
\end{solution}

% % Question 3
% \question[2]
%     If $f(x)$ is an increasing function, determine if $g(x) = f(-x)$ is increasing, decreasing, or neither, and justify.

% % Solution
% \begin{solution}[0cm]
%     $g(x)$ is a decreasing function.
% \end{solution}

% Question 4
\question[4]
    Consider the function $f(x) = x^3 - 3x^2 + 4$ defined on the interval $x \in [0, 3]$
    \begin{enumerate}[leftmargin = 0.5cm, label = (\alph*), rightmargin = 0.2cm]
        \item Compute $f'(x)$.
        \item Find the critical points of $f(x)$ in the interval $[0, 3]$.
        \item Find the local extreme points of $f(x)$.
        \item Find the global maximum and minimum of $f(x)$ on the interval $[0, 3]$.
        \item Do the global extremes change if the domain is taken to be $\mathbb{R}$ instead of $[0, 3]$? Why?
    \end{enumerate}

% Solution
\begin{solution}[0cm]
    \begin{enumerate}[leftmargin = 0.5cm, label = (\alph*), rightmargin = 0.2cm]
        \item $f'(x) = 3x^2 - 6x = 3x(x - 2)$.
        \item The critical points in $[0, 3]$ are $x = 0$ and $x = 2$.
        \item There is a local maximum at $x = 0$ and a local minimum at $x = 2$.
        \item $f(0) = 4$, $f(2) = 0$, and $f(3) = 4$, therefore the global maximum is $4$ (attained at $x = 0$ and $x = 3$) and the global minimum is $0$ (attained at $x = 2$)
        
        \item $\lim_{x\rightarrow\infty} f(x) = \infty$ and $\lim_{x\rightarrow-\infty} f(x) = -\infty$, therefore there is no global maximum or minimum when the domain is considered to be $\mathbb{R}$.
    \end{enumerate}
\end{solution}

\end{questions}	

\begin{center}
You may find useful the following derivatives of basic functions and rules of derivatives
\end{center}

\begin{center}
    \textbf{Differentiation rules}
\end{center}
        
        \begin{enumerate}[leftmargin = 0.5cm, label = , rightmargin = 0.2cm]

        \item \textbf{Linearity:} $h(x) = a f(x) + b g(x) \Rightarrow h'(x) = a f'(x) + b g'(x)$, for $a, b \in\R$

        \item \textbf{Product rule:} $h(x) = f(x)g(x) \Rightarrow h'(x) = f'(x)g(x) + f(x)g'(x)$

        \item \textbf{Quotient rule:} $h(x) = f(x)/g(x) \text{ and } g'(x) \neq 0 \Rightarrow h'(x) = \frac{f'(x)g(x) - f(x)g'(x)}{g(x)^2}$

        \item \textbf{Chain rule:} $h(x) = g(f(x)) \Rightarrow h'(x) = g'(f(x))f'(x)$

        \end{enumerate}

\begin{center}
    \textbf{Derivatives of basic functions}
\end{center}

        \begin{enumerate}[leftmargin = 0.5cm, label = , rightmargin = 0.2cm]

        \item \textbf{Constant:} $f(x) = a \Rightarrow f'(x) = 0$
        
        \item \textbf{Power:} $f(x) = x^a \Rightarrow f'(x) = a x^{a-1}$
        
        \item \textbf{Exponential:} $f(x) = e^x \Rightarrow f'(x) = e^x$

        \item \textbf{Logarithm:} $f(x) = \ln(x) \Rightarrow f'(x) = 1/x$
        
        \end{enumerate}

\end{document}
